\documentclass[11pt]{article}
\usepackage[a4paper,margin=26mm]{geometry}
\usepackage[T1]{fontenc}
\usepackage{lmodern,amsmath,amssymb,amsthm,booktabs,microtype}
\usepackage[colorlinks=true,allcolors=blue]{hyperref}
\newtheorem{theorem}{Theorem}
\newtheorem{lemma}[theorem]{Lemma}
\newtheorem{proposition}[theorem]{Proposition}
\newtheorem{corollary}[theorem]{Corollary}
\newcommand{\eps}{\varepsilon}
\newcommand{\tr}{\operatorname{tr}}
\newcommand{\adj}{\operatorname{adj}}
\newcommand{\Prob}{\mathbb P}
\newcommand{\E}{\mathbb E}
\title{An exact two-mean law for singleton--avoidance probabilities}
\author{Lluis Eriksson}
\date{24 October 2026 --- research draft}
\begin{document}
\maketitle
\begin{abstract}
We determine the infimum of $\Prob(X=1,Z=0)$ over finitely supported
nonnegative integer-valued pairs with real-stable probability generating
polynomial and prescribed means $(\alpha,b)$. In the positive region
$0<\alpha<2$, $0\le b<\min(1,2-\alpha)$, the answer is the smaller of
a degree-two value and an explicit categorical--Poisson limit. For
$\alpha\le1$ the limiting categorical trial has no $X$ mass; for
$\alpha>1$ and $b>0$ it necessarily has positive $X$ mass. A determinantal
pinching preserves the coefficient and both means exactly. Three
perturbations of a hypothetical finite global minimizer then exclude
all intermediate Bernoulli block sizes, using an analytic estimate for
large blocks and three elementary polynomial factorizations for the
remaining sizes. The bound is sharp even for finite determinantal measures
with Hermitian positive-contraction kernels, counted on two disjoint
blocks. The formula concerns unrestricted finite
degree and one singleton target; no application to a TSP approximation
constant is asserted.
\end{abstract}

\paragraph{Draft status.}
The preceding exact draft received a single-model assessment recommending
minor revision. This corrected draft awaits a check of those changes and
is not published. Exact algebraic checks supplement the proofs below.

\section{The exact envelope}
A real polynomial is \emph{real stable} if it is nonzero whenever all
variables have positive imaginary parts. Let $\mathcal S(\alpha,b)$
consist of real-stable polynomials with nonnegative coefficients,
$P(1,1)=1$, and $\nabla P(1,1)=(\alpha,b)$. Degrees are finite but have
no common bound. Define
\[
 C(\alpha,b)=\inf_{P\in\mathcal S(\alpha,b)}[xz^0]P.
\]
These classes are nonempty for all nonnegative means: independent
monomials times Bernoulli factors realize every such pair.
For the positive region put
\begin{align}
 \eps&=1-b,& U&=\min(\alpha,\eps),\\
 a_1&=\min\{U,\max\{0,(2\alpha-b)/4\}\},&
 H&=\eps\alpha-(2\alpha-b)a_1+2a_1^2.\label{eq:H}
\end{align}
Define $a_\infty=0$ when $\alpha\le1$. For $1<\alpha<2$, set
\begin{equation}\label{eq:ainfty}
 a_\infty=\frac{2b(\alpha-1)}{2-\alpha+b+
 \sqrt{(2-\alpha+b)^2+4b(\alpha-1)}}.
\end{equation}
This includes $b=0$. Finally put
\begin{equation}\label{eq:J}
 J=e^{a_\infty-\alpha}
 \{\eps(\alpha-a_\infty)+a_\infty(1-\alpha+a_\infty)\}.
\end{equation}
\begin{theorem}\label{thm:main}
For $\alpha,b\ge0$, the envelope is zero if $\alpha=0$, $b\ge1$,
$\alpha\ge2$, or $\alpha+b\ge2$. At every remaining point,
\begin{equation}\label{eq:main}
 C(\alpha,b)=\min\{H,J\}>0.
\end{equation}
The value $H$ is attained by a polynomial of degree at most two; $J$ is
approached by normalized categorical--binomial products. The same
infimum holds for disjoint-block counts of strongly Rayleigh measures,
and even when restricted to finite determinantal measures with Hermitian
positive-contraction kernels $0\preceq K\preceq I$.
\end{theorem}
The zero cases in the theorem overlap harmlessly. The positive domain
is precisely $0<\alpha<2$, $0\le b<\min(1,2-\alpha)$.
When $\alpha>1$, the quadratic vertex satisfies
$(2\alpha-b)/4-(\alpha-1)=(4-2\alpha-b)/4>0$ in this domain,
so the definition of $a_1$ also respects the lower feasibility bound
$a\ge\alpha-1$ for the degree-two realizer.
This domain was already covered by refined positivity estimates in
Gurvits--Klein--Leake \cite[Theorem 5.11 and Corollary 5.13]{GKL}.
Our question is the exact value, including the competing finite and
limiting mechanisms. We do not claim a new positivity region.

For $\alpha\le1$, the limiting branch simplifies to
$J=\alpha(1-b)e^{-\alpha}$. For $\alpha>1$, $b>0$,
$a_\infty\in(0,\alpha-1)$: using a pure Poisson block and a disjoint
avoidance trial misses the smaller mixed limit. An exact finite example
illustrating that failure is given in Section~\ref{sec:example}.

The deep input is Borcea--Br\"and\'en's bivariate determinantal
representation \cite[Corollary 6.7]{BB}. The Bernoulli equalization
argument is classical Hoeffding theory \cite{Hoeffding}; we give the
needed proof. Capacity productization in \cite{GL,GKL} does not by
itself assert exact preservation of our boundary coefficient.
These comparisons identify dependencies and scope, not worldwide
originality. Higher targets and several simultaneous positive target
counts remain outside the theorem.

\section{Exact coefficient preservation}
\begin{proposition}\label{prop:pinch}
For $P\in\mathcal S(\alpha,b)$ of exact degree $D\ge1$, there is a product
$Q=\prod_{i=1}^D(c_i+a_ix+b_iz)$ with nonnegative entries,
$a_i+b_i+c_i=1$, and $a_i+b_i>0$, having the same means and coefficient
$[xz^0]$ as $P$.
\end{proposition}
\begin{proof}
The representation in \cite{BB} gives
$P=\pm\det(xA+zB+C)$ with real symmetric $D$ by $D$ matrices,
$A,B\succeq0$. The coefficient of $t^D$ in $P(t,t)$ is positive and
is $\pm\det(A+B)$. Thus $G=A+B\succ0$ and the sign is positive.
Since
\[
 P(t,t)=\det G\det(tI+G^{-1/2}CG^{-1/2})>0\quad(t>0),
\]
the symmetric matrix $G^{-1/2}CG^{-1/2}$ has no negative eigenvalue.
Hence $C\succeq0$. Now $A+B+C\succ0$ has determinant one; congruence
by its inverse square root preserves $P$ and normalizes the sum to $I$.
Jacobi's polynomial identity yields
\[
 \partial_xP(1,1)=\tr A,\quad \partial_zP(1,1)=\tr B,
 \quad [xz^0]P=\tr(\adj(C)A).
\]
The last identity also holds for singular $C$. Diagonalize $C$ and
replace $A,B$ by their diagonals in that basis. This preserves both
traces and $\tr(\adj(C)A)$, since $\adj(C)$ is diagonal. The diagonal
entries are the required $a_i,b_i,c_i$. Their sums are one, and
$A+B\succ0$ implies $a_i+b_i>0$. The determinant becomes $Q$.
\end{proof}
No inverse of $C$ is used. The proposition preserves a coefficient and
means, not the full distribution or higher coefficients. The converse
inclusion holds because the displayed products themselves are stable.

\section{Reduction to one categorical trial and an equal block}
Throughout this section $0<b<\min(1,2-\alpha)$ and $0<\alpha<2$.
For $k\ge2$, $0\le a\le U$, put
\begin{equation}\label{eq:G}
 G_k(a)=a(1-u)^k+(\eps-a)(\alpha-a)(1-u)^{k-1},
 \qquad u=(\alpha-a)/k.
\end{equation}
For $k=1$, use the same polynomial formula on
$\max(0,\alpha-1)\le a\le U$. These values are realized by
\begin{equation}\label{eq:realizer}
 (ax+bz+\eps-a)(1-u+ux)^k.
\end{equation}
A factor with $u=0$ can be dropped.

\begin{lemma}\label{lem:equal}
Fix a finite number of Bernoulli slots, their total mean $m<2$, and
nonnegative weights $A,C$. A minimum of
$A\Prob(S=0)+C\Prob(S=1)$ can be chosen with all fractional Bernoulli
parameters equal and at most one parameter equal to one.
\end{lemma}
\begin{proof}
The fixed-slot feasible set is compact. Among its minimizers choose one
with fewest fractional parameters. For two such parameters $u,v$ with
fixed sum, the objective is affine in $uv$. A positive coefficient
forces an endpoint, a zero coefficient permits an endpoint with fewer
fractional entries, and a negative coefficient forces $u=v$. Thus all
fractional entries are equal. There is at most one deterministic one
because the total mean is less than two. Zero entries may be removed.
\end{proof}
The fixed number of slots is chosen \emph{before} applying compactness;
only afterwards do we take the infimum over all finite slot counts.

\begin{proposition}\label{prop:scalar}
$C(\alpha,b)$ equals the infimum of all the values in \eqref{eq:G},
including the stated $k=1$ domain.
\end{proposition}
\begin{proof}
Use Proposition~\ref{prop:pinch}. At fixed constants $c_i$, the target
coefficient $\sum_i a_i\prod_{j\ne i}c_j$ is linear on
$0\le a_i\le1-c_i$, $\sum_i a_i=\alpha$. A vertex has at most one
coordinate strictly between its bounds. Therefore a minimizing
replacement has at most one factor containing both $x$ and $z$.
This reasoning remains valid when some constants vanish.

If there is a mixed factor, write $a$ for its $x$ mass and $t$ for the
sum of pure $z$ masses. Merging the latter into $1-t+tz$ decreases the
coefficient, since their constant product is at least $1-t$. The mixed
factor becomes $ax+(b-t)z+\eps-a+t$. For the remaining pure $x$ block
write $P_0,P_1$ for its zero and one probabilities; its mean is
$\alpha-a<2$. The objective is
\[
 (1-t)\{aP_0+(\eps-a+t)P_1\},\qquad
 \max(0,a-\eps)\le t\le b.
\]
This is concave in $t$, so an endpoint is no larger.
At $t=0$, $a\le U$, and Lemma~\ref{lem:equal} gives an equal block,
except possibly for a deterministic one. In that case $\alpha\ge1$,
$a\le\alpha-1$, and the objective is at least
\[
 (\eps-a)(2-\alpha+a)\ge\delta,
 \qquad \delta=2-\alpha-b.
\]
Indeed the left side is concave on $[0,\alpha-1]$ and its endpoint
values are $\eps(2-\alpha)$ and $\delta$; their difference is
$b(\alpha-1)\ge0$. The value $\delta=G_1(\alpha-1)$ is feasible.
The empty block has $a=\alpha\le\eps$ and is included with $u=0$.

At $t=b$, or with no mixed factor to start with, the objective is
$\eps$ times a pure Bernoulli one-probability of mean $\alpha$.
Equal blocks give $G_k(0)$; the deterministic alternative is at least
$\eps(2-\alpha)\ge\delta$ when it exists.
At the other possible endpoint $t=a-\eps>0$ one has
$a>\eps>\alpha-1$, so $\alpha-a<1$. Then
\begin{align*}
 a(1-a+\eps)P_0&\ge a(1-a+\eps)(1-\alpha+a)\\
 &\ge\eps(1-\alpha+\eps)=G_1(\eps).
\end{align*}
For the second inequality subtract the right side and factor:
\[
 (a-\eps)\{(1-a)(1-\alpha+a)+\eps\}\ge0.
\]
Here $a\le1$, and $\eps<a\le\alpha$ makes $G_1(\eps)$ feasible.
All reductions are thus dominated by the scalar union. Conversely,
\eqref{eq:realizer} realizes every scalar value.
\end{proof}

\section{The limit and the two boundaries}
Uniformly on $a\in[0,U]$, $G_k(a)$ tends to
\begin{equation}\label{eq:L}
 L(a)=e^{a-\alpha}\{\eps(\alpha-a)+a(1-\alpha+a)\}.
\end{equation}
For example, for $k\ge4$ and $0\le\mu=\alpha-a<2$,
\[
 |k\log(1-\mu/k)+\mu|\le
 \frac{\mu^2}{2k(1-\mu/k)}\le4/k.
\]
Together with \eqref{eq:G} this proves uniform convergence, including
when $a$ varies with $k$. The sign of $L'(a)$ is the sign of
\[
 a^2+(2-\alpha+b)a-b(\alpha-1).
\]
It is increasing on $a\ge0$. Its minimum occurs at zero for
$\alpha\le1$, and otherwise at \eqref{eq:ainfty}, strictly between
zero and $\alpha-1$. Thus $\min L=J$. Also $\min G_1=H$.

We will exclude any global scalar minimum below both $H$ and $J$.
Uniform convergence ensures that such an infimum is attained in a
finite collection of compact scalar families: the large-$k$ tail is
separated from it by a positive gap.
At $a=U=\alpha$ the value is $\alpha=G_1(\alpha)\ge H$.
At $a=U=\eps$, $\mu=\alpha-\eps=1-\delta\in[0,1)$, and
\[
 G_k(\eps)=\eps(1-\mu/k)^k\ge\eps(1-\mu)
 =\eps\delta=G_1(\eps)\ge H.
\]
To handle $a=0$, we record the classical univariate boundary explicitly.

\begin{lemma}\label{lem:uni}
For a finite sum of independent Bernoulli variables of mean $\alpha$,
$0<\alpha<2$, the infimum of its one-probability is
\[
 M(\alpha)=\begin{cases}
 \alpha e^{-\alpha},&0<\alpha\le1,\\
 \min\{\alpha e^{-\alpha},\alpha(1-\alpha/2)\},&1<\alpha<2.
 \end{cases}
\]
\end{lemma}
\begin{proof}
Lemma~\ref{lem:equal} reduces the objective to equal blocks or one
deterministic success. The latter gives at least $2-\alpha$, which
exceeds $\alpha(1-\alpha/2)$ by $(2-\alpha)^2/2$.
Equal blocks have $q(k)=\alpha(1-\alpha/k)^{k-1}$.
For $\alpha\le1$, $q(k)>\alpha e^{-\alpha}$ for every finite feasible
$k$, by $\log(1-t)>-t/(1-t)$; the deterministic $k=1$ endpoint is
separate. For $\alpha>1$, extend $k$ to real $k\ge2$ and put
$t=\alpha/k$. Direct differentiation gives
\[
 (1-t)\frac{d\log q}{dk}=t^2\{R(t)-1/\alpha\},\qquad
 R(t)=\frac{t+(1-t)\log(1-t)}{t^2}
 =\sum_{j\ge0}\frac{t^j}{(j+1)(j+2)}.
\]
$R$ increases with $t$, so the derivative sign decreases as $k$
increases. A critical point can only be a maximum. The infimum is
therefore at $k=2$ or $k\to\infty$, as stated.
\end{proof}
At $a=0$, Lemma~\ref{lem:uni} excludes a value below $\min(H,J)$.
For $\alpha\le1$ its lower bound is $J$.
For $\alpha>1$ the Poisson bound is $L(0)\ge J$, while the quadratic
bound $\eps\alpha(1-\alpha/2)$ is at least $H$. To check the latter,
put $d=2-\alpha$, $\delta=d-b$. The unconstrained quadratic branch is
\[
 Q=\delta-(b+2\delta)^2/8=\alpha\delta/2-b^2/8;
\]
$H=Q$ if $b\le2\delta$, and $H=E:=\eps\delta$ otherwise.
The differences from these two branches are respectively
$b\alpha(\alpha-1)/2+b^2/8>0$ and
$\eps(b-d^2/2)>0$; in the clipped branch $b\ge2d/3>d^2/2$.

\section{Three variations exclude intermediate finite blocks}
Assume a scalar value lies below $\min(H,J)$, and select an attained
global minimum $V$ as above. It has $k\ge2$, $0<a<U$, and
$p=(\alpha-a)/k\in(0,1)$. Lemma~\ref{lem:equal} and the deterministic
bound in Proposition~\ref{prop:scalar} show that every finite admissible
product with the fixed means $(\alpha,b)$ has value at least $V$.
This applies also after changing $a$ or appending a zero Bernoulli slot;
both types of perturbation below are therefore covered.
Write $\mu=kp$, $q=\eps-a$, $y=(k+1)p$, and $B=2-y$.

First change $a$ to $a+t$ and one pure probability $p$ to $p-t$,
keeping $b$ fixed. This preserves the two means and is feasible for
small $t$ of either sign. Dividing its linear and quadratic
coefficients by the positive factor $(1-p)^{k-2}$ gives
\[
 S=B(a-p)-b(\mu-1)=0,\qquad B\ge0.
\]
If $B=0$, then $\mu=2-p>1$ and $S<0$, impossible. Set
$\beta=b/B>0$. Stationarity gives
\begin{align}
 a&=p+\beta(\mu-1),& q&=(1-p)(1-\beta),\nonumber\\
 \alpha&=y+\beta(\mu-1),&\delta&=B-\beta(1-p),\label{eq:param}\\
 G_k&=(1-p)^k(y-\beta),&b&=\beta B.\nonumber
\end{align}
Second perturb two pure probabilities to $p+t,p-t$. The coefficient
of their product is a positive factor times
$a(1-p)+q(\mu-2)=(1-p)(\beta-B)$. Since their product decreases
by $t^2$, minimality requires $\beta\le B$.
Third split one probability into $p-t,t$, $0<t<p$, using a new slot.
The exact objective change is
\[
 t(p-t)(1-p)^{k-2}\{a(1-p)+q(y-2)\}.
\]
The bracket equals $(1-p)(p-\delta)$ by stationarity, hence
$\delta\le p$. Combining both necessary conditions gives
\[
 B-p\le\beta(1-p)\le B(1-p),\qquad B\le1.
\]
Equality $B=1$ forces $\beta=1$ and $b=1$, excluded. Thus
\begin{equation}\label{eq:domain}
 0<B<1,\quad n=k+1\ge3,\quad p=(2-B)/n,\quad
 \max(0,\ell)\le\beta\le B,\quad
 \ell=\frac{(n+1)B-2}{n-2+B}.
\end{equation}
This deduction does not require $\alpha>1$.

\subsection{The quadratic branch remains feasible}
Before comparing with $H$, its branch formula must be justified also
when the original mean is below one. At a stationary point,
\[
 2\alpha-b=4-2B+\beta(2-3B-2p)>0.
\]
Indeed, if the coefficient of $\beta$ is negative, use $\beta\le B$
and $p\le(2-B)/3$ to bound it below by
$4-7B^2/3-4B/3>1/3$; the other sign is immediate.
Thus $a_*=(2\alpha-b)/4$ is positive and at most $\alpha$.
Furthermore $a_*-(\alpha-1)=(4-2\alpha-b)/4>0$.
Only the constraint $a_*\le\eps$ may clip it. Consequently at these
points
\begin{equation}\label{eq:branches}
 H=\begin{cases}Q,&b\le2\delta,\\E,&b\ge2\delta,\end{cases}
 \quad Q=\delta-(b+2\delta)^2/8,\quad E=(1-b)\delta.
\end{equation}
The identity $Q-E=-(b-2\delta)^2/8\le0$ also proves $H\le E$
even if the formal endpoint $a=\eps$ would
not have been feasible before imposing the stationary conditions.

\subsection{All $k\ge5$}
The ratio $(y-\beta)/(1-\beta B)$ decreases with $\beta$; its
derivative is $-(y-1)^2/(1-\beta B)^2$. Hence
\[
 \frac{G_k}{\eps}\ge A_n(y):=
 \frac{2(1-y/n)^{n-1}}{3-y},\qquad 1<y<2.
\]
For $n\ge6$,
\begin{equation}\label{eq:large}
 A_n(y)\ge\min(y/n,2-y).
\end{equation}
To prove this, put $y_*=2n/(n+1)$. On $[1,y_*]$ the logarithmic
derivative of $A_n(y)/(y/n)$ is
$-(n-1)/(n-y)+1/(3-y)-1/y<0$.
On $[y_*,2)$ that of $A_n(y)/(2-y)$ is
$-(n-1)/(n-y)+1/(3-y)+1/(2-y)>0$:
here $y_*\ge3/2$, $1/(2-y)\ge2$, and $(n-1)/(n-y)\le2$.
The minimum ratio is at $y_*$ and equals
\[
 R_n=\frac{(n+1)^2}{n+3}
 \left(\frac{n-1}{n+1}\right)^{n-1}.
\]
$R_6=3125/3087>1$, and for real $n\ge6$,
\[
 \frac{d\log R_n}{dn}
 \ge\frac4{n+1}-\frac1{n+3}-\frac2{n-1}
 =\frac{n^2-17}{(n-1)(n+1)(n+3)}>0,
\]
using $\log(1-u)\ge-u/(1-u)$. This proves \eqref{eq:large}.
Since $\delta\le\min(p,B)$, it gives $G_k\ge\eps\delta=E\ge H$,
contradicting the violating minimum.

\subsection{The three remaining sizes}
Under \eqref{eq:domain}, the exact factorizations in
Appendix~\ref{app:small} prove $G_k\ge H$ for $n=3,4,5$.
Thus no finite global minimum can lie below both $H$ and $J$.
The uniform-convergence argument proves the lower bound everywhere.
Conversely \eqref{eq:realizer} with $k=1,a=a_1$ attains $H$ and
with $a=a_\infty,k\to\infty$ approaches $J$.
This proves \eqref{eq:main} for $b>0$.

\section{Zero region, boundary and determinantal sharpness}
At $b=0$, nonnegative coefficients force $P$ to be univariate in $x$.
Its real nonpositive roots factor it into normalized Bernoulli factors;
Lemma~\ref{lem:uni} proves exactly \eqref{eq:main}.
At $\alpha=0$ the target coefficient vanishes. If $b\ge1$, a factor
$z$ ensures its vanishing, with remaining means realized independently.
If $\alpha\ge2$, use a factor $x^2$. For $1<\alpha<2$ and
$b\ge2-\alpha$, use
\[
 \bigl((\alpha/2)x+(1-\alpha/2)z\bigr)^2
\]
and an independent $z$-only count of mean $b-(2-\alpha)$.
These cover the remaining zero cases. Positivity inside the stated
region follows also from the formula: each feasible degree-two product
has positive target value, and the limiting prefactor in \eqref{eq:J}
is positive.

Here a finite determinantal measure means one with a Hermitian kernel
$0\preceq K\preceq I$; such measures are strongly Rayleigh
\cite[Proposition 3.5]{BBL}.
If a strongly Rayleigh measure counts disjoint blocks $A,B$, identifying
the variables in each block and setting the other variables to one
produces a bivariate stable polynomial with the same means and event.
Conversely every categorical factor $c+ax+bz$ is realized on a two-point
block by the determinantal kernel
\[
 K=\begin{pmatrix}a&\sqrt{ab}\\\sqrt{ab}&b\end{pmatrix},
 \qquad 0\preceq K\preceq I.
\]
Its probabilities are $c,a,b,0$ for neither, the first point, the second
point and both. Independent blocks realize the products in
\eqref{eq:realizer}; one-point Bernoulli blocks suffice for pure factors.
The upper constructions therefore belong to the determinantal class,
which itself lies in the strongly Rayleigh class \cite{BBL}.
This proves the sharpness assertion of Theorem~\ref{thm:main}.
Additional avoided blocks can all be identified with $z$; only the sum
of their expected counts is then prescribed.

\section{A mixed limit and the boundary of the result}\label{sec:example}
At $\alpha=13/10,b=3/10$, the pure-Poisson candidate is
$(91/100)e^{-13/10}$ and the degree-two value is $199/800$.
The valid finite product
\[
 \left(\frac{x}{10}+\frac{3z}{10}+\frac35\right)
 \left(\frac{247}{250}+\frac{3x}{250}\right)^{100}
\]
has coefficient
\[
 \left(\frac{247}{250}\right)^{99}
 \left(\frac{41}{50}-\frac3{2500}\right)
 =0.2478080145\ldots,
\]
strictly below both naive candidates. The exponential comparison can
be certified by rational Taylor bounds; this finite example is not
used as proof of the main theorem. The correct exceptional mass is
$a_\infty=(\sqrt{34}-5)/10$.
A further tempting shortcut would compare $G_k(a)$ to $G_1(a)$ and
$L(a)$ at the \emph{same} $a$; that assertion is false. For
$\alpha=1,b=13/20,a=1/5,k=3$,
$G_3=484/3375$ is smaller than both $G_1=4/25$ and
$L=(8/25)e^{-4/5}$. The global quadratic minimum, however, is
$391/3200<G_3$. The proof therefore requires a global minimizer and
feasible perturbations, rather than a pointwise block-size comparison.

The theorem does not give a fixed-degree envelope, classify all equality
cases, preserve higher coefficients under pinching, or handle several
positive targets simultaneously. Its limiting branch denotes an
infimum approached by finite polynomials, not a finite polynomial
identified with a Poisson law. The univariate ingredients are classical;
the determinantal representation is a cited deep theorem. Neither a
new TSP approximation ratio nor exhaustive priority is asserted.

\paragraph{AI and responsibility disclosure.}
The named author directed substantial OpenAI Codex assistance in
derivation, writing and exact algebraic work. GPT-6 Astra xhigh audits
used shared files and declared reused context. These checks are not
independent human peer review or formal proof-assistant verification.
The exact preceding draft and its single-model assessment are retained
separately; this revision does not itself constitute a new full assessment.
The author is responsible for the submitted text and claims.

\appendix
\section{Three exact small-block comparisons}\label{app:small}
Here $n=3,4,5$, $p=(2-B)/n$, and the constraints are
\eqref{eq:domain}. Write $t=\beta$ and
\[
 G=(1-p)^{n-1}(2-B-t),\quad \delta=B-t(1-p),\quad b=tB.
\]
Let $Q,E$ be as in \eqref{eq:branches}. The branch threshold is
\[
 T=\frac{2B}{B+2(1-p)};
 \qquad H=Q\ (t\le T),\quad H=E\ (t\ge T).
\]
The identity $G-Q=(2-B)^2P_n/D_n$ holds with
\begin{align*}
 P_3&=t^2+4t+4-8B,&D_3&=72,\\
 P_4&=2t^2+(6-B)t+4-8B-B^2,&D_4&=64,\\
 P_5&=225t^2+(588-128B-8B^2)t\\
 &\hspace{8mm}+324-656B-112B^2-8B^3,&D_5&=5000.
\end{align*}
Each $P_n$ increases with $t\ge0$, $0<B<1$. At $t=0$ and
$B\le2/(n+1)$ its nonconstant-in-$B$ expression is nonnegative:
for $n=3$ it is $4-8B$; for $n=4$ its decreasing value at $2/5$
is $16/25$; for $n=5$ its decreasing value at $1/3$ is $2500/27$.
For $B\ge2/(n+1)$, evaluate at $t=\ell$:
\begin{align*}
 (G-Q)_{n=3}&=\frac{B(2-B)^3(2B-1)}{18(B+1)^2},\\
 (G-Q)_{n=4}&=\frac{B(2-B)^3(B^2+19B-2)}{64(B+2)^2},\\
 (G-Q)_{n=5}&=\frac{(2-B)^3(B^4+28B^3+343B^2+60B+18)}{625(B+3)^2}.
\end{align*}
These are nonnegative on the respective ranges. Thus $G\ge Q$
throughout the relaxed feasible set.

For the clipped branch, $G-E$ is concave in $t$, with second derivative
$-2B(1-p)<0$. Its feasible interval is $[\max(0,\ell,T),B]$.
If nonempty, $T\le B$ implies $B\ge B_c=4/(n+2)$.
At $t=B$ its values are
\begin{align*}
 n=3:&\quad(1-B)(B+1)(3B^2-4B+2)/9,\\
 n=4:&\quad(1-B)(9B^3-2B^2-4B+8)/32,\\
 n=5:&\quad(1-B)(2B^4+149B^3-17B^2-34B+162)/625.
\end{align*}
They are positive: complete the square in the first, and bound the
negative terms by $6$ and $51$ respectively in the last two.
At the lower endpoint $T$, $Q=E$ and the preceding comparison applies.
If the lower endpoint is instead $\ell>T$, the exact values are
\[
 \frac{(1-B)R_n(B)}{d_n(B+n-2)},\qquad
 (d_3,d_4,d_5)=(9,64,625),
\]
where
\begin{align*}
 R_3&=B^3+18B^2-12B-2,\\
 R_4&=B^4+12B^3+128B^2-48B-16,\\
 R_5&=B^5+20B^4+150B^3+1290B^2-180B-102.
\end{align*}
Each increases for $B\ge B_c$: the positive linear term in its
derivative already dominates the sole negative constant.
Their values at $B_c$ are $54/125$, $1024/81$, and
$4143750/16807$, all positive. Concavity now gives $G\ge E$
on the clipped interval. This proves all three remaining comparisons.
The displayed identities are rational polynomial identities; their
checks require no numerical minimization or sampled sign inference.

\begin{thebibliography}{9}
\bibitem{BB} J. Borcea and P. Br\"and\'en.
Multivariate P\'olya--Schur classification problems in the Weyl algebra.
\emph{Proc. London Math. Soc.} 101 (2010), 73--104.
\href{https://arxiv.org/abs/math/0606360v6}{arXiv:math/0606360v6}, Corollary 6.7.
\bibitem{BBL} J. Borcea, P. Br\"and\'en and T. M. Liggett.
Negative dependence and the geometry of polynomials.
\emph{J. Amer. Math. Soc.} 22 (2009), 521--567.
\href{https://arxiv.org/abs/0707.2340}{arXiv:0707.2340}.
\bibitem{Hoeffding} W. Hoeffding.
On the distribution of the number of successes in independent trials.
\emph{Ann. Math. Statist.} 27 (1956), 713--721.
\href{https://www.jstor.org/stable/2237378}{Journal record}.
\bibitem{GL} L. Gurvits and J. Leake.
Capacity lower bounds via productization. \emph{STOC 2021}, 847--858.
\href{https://eccc.weizmann.ac.il/report/2020/110/}{ECCC TR20-110}.
\bibitem{GKL} L. Gurvits, N. Klein and J. Leake.
From trees to polynomials and back again: New capacity bounds with
applications to TSP. \emph{ICALP 2024}, 79:1--79:20.
\href{https://arxiv.org/abs/2311.09072v2}{arXiv:2311.09072v2}.
\end{thebibliography}
\end{document}
